\chapter{いろいろな素材と便利な機能}\label{chap:sources}

% ------------------------------------------------------------------------------
\section{月・太陽を処理する}\label{sec:moon}

月面や太陽面は、画面いっぱいに模様が広がっています。処理の手順は惑星と同じで、多くの場合は既定値のままで大丈夫です。

\screen[0.78\linewidth]{moon_ap}{月面の参照画像と位置合わせ領域（1024×768、自動で 128 画素・124 個）}{}

\begin{itemize}[itemsep=2pt]
  \item \textbf{対象モード}は\uil{自動}で\uil{月・太陽}と判定されます。うまくいかないときはアライメントタブで\uil{月・太陽}を選びます。
  \item 位置合わせ領域は、画面の明るい所いっぱいに並びます。上の例では欠け際の暗い所には置かれていません。
        暗い所にも置きたいときは、詳細設定の\uil{配置する明るさ}を下げます（第\ref{chap:align}章）。
  \item 画面が大きい（例：1920×1200 以上）と時間がかかります。対象が画面の一部だけなら\textbf{処理範囲}（第\ref{sec:roi-howto}節）を使うと速くなります。
  \item ウェーブレットは、Layer 1〜2 を控えめ（2〜5 くらい）、Layer 3〜4 を少し上げると自然な仕上がりになりやすいです。
  \item 太陽の Hα など 16bit のモノクロ SER では、画面上部に「ヘッダと異なるバイトオーダーで読んでいます」と灰色で出ることがあります。
        画像が正常に見えていれば問題ありません（第\ref{sec:byteorder}節）。
\end{itemize}

% ------------------------------------------------------------------------------
\section{カメラの RAW・静止画を処理する}\label{sec:raw}

一眼レフ・ミラーレスで連写した RAW（CR2・CR3・NEF・ARW・RAF など）や、TIFF・PNG・JPEG・FITS の連番も扱えます。

\begin{enumerate}[itemsep=2pt]
  \item 画像を1つのフォルダにまとめ、\textbf{フォルダごと}入力キューへドロップします（複数の画像を選んでドロップしてもかまいません）。
        名前の順に並べた1本の「連番」として追加されます。
  \item あとは動画と同じく \ui{品質評価}\ui{アライメント}\ui{スタック} と進めます。
  \item RAW は、カメラの色の調整（ホワイトバランス・色の変換・明るさの曲線）を掛けない\textbf{そのままの明るさ（リニア）}で読みます。
        そのため、スタックの結果は\textbf{緑がかって暗く}見えます。仕上げの \ui{自動（灰色仮説）} で色を、レベル補正で明るさを整えてください。
\end{enumerate}

\begin{tip}[写真の枚数が少ないとき]
  連写が数枚〜数十枚しかないときは、スタックタブの採用率を大きく（30〜100\%）、アライメントタブの参照フレームも大きく（50\%）すると、
  ノイズの少ない結果になります。
\end{tip}

% ------------------------------------------------------------------------------
\section{処理範囲を使う}\label{sec:roi-howto}

大きなセンサー（例：6240×4160 画素）で撮り、月や惑星が画面の中央に小さく写っているだけなら、\textbf{処理範囲}を決めてその周りだけを処理すると、
何倍も速くなり、メモリも少なくて済みます。

\screen{cr2_roi}{カメラの RAW（8枚）で処理範囲を決めたところ}{%
  \mk{queuerow}{1}\mkd[nw]{roibox}{2}\mk[w]{draw}{3}\mk[w]{label}{4}\mk[w]{clear}{5}\mk[w]{status}{6}
  \drag{($(roibox-nw)+(0.015,-0.03)$)}{($(roibox-se)+(-0.015,0.03)$)}}

\begin{marklist}
  \item フォルダを追加すると、\uil{cr2（静止画連番）・8フレーム・6240×4160} のように1本の連番として並びます。
  \item \textbf{処理範囲}（黄色の点線）。品質評価タブの \uil{フレームの上で処理範囲を描く} をオンにし、プレビューの上でドラッグして描きます。
        枠の外でドラッグすると描き直し、内側で移動、辺や角で大きさを変えられます。
  \item 範囲を描くモードのオン・オフ。
  \item 決まった範囲（例：3120×2080 画素、左上が x 1560・y 1040）。
  \item \ui{全体に戻す}：範囲を消して全体を処理します。
  \item 範囲を変えると、品質評価からやり直しになります。
\end{marklist}

\begin{caution}[対象が範囲からはみ出さないように]
  撮影中、対象は少しずつ動きます。\textbf{余白を取って}囲み、フレームのスライダーを最後まで送って、対象が範囲に収まり続けているか確かめてください。
  品質評価・アライメントのあとに、はみ出すフレームがあればその数を知らせます。
\end{caution}

% ------------------------------------------------------------------------------
\section{MOV・MP4 の動画を処理する}\label{sec:movie}

スマートフォンやデジタルカメラの MOV・MP4・M4V（H.264・HEVC・ProRes など）も開けます。
HEVC（H.265）の動画は、macOS 10.13 以降でだけ読めます（10.12 には macOS 側に HEVC の読み込み機能がありません）。

\screen{mp4_banner}{MOV・MP4 を開いたときの注意の帯}{\mk[inw]{banner}{1}}

\begin{marklist}
  \item MOV・MP4 は macOS の仕組み（デコーダ）で読むため、\textbf{macOS の版によって画素の値がわずかに変わることがあります}。
        同じ結果を何度でも再現したい記録用途には、撮影ソフトで SER や AVI（非圧縮）で撮るのがおすすめです。\ui{×} で閉じられます。
\end{marklist}

\begin{caution}
  圧縮された動画は、細かい模様が圧縮で失われていることがあります。また、縦向きで撮ったスマートフォンの動画は横向きに読まれます
  （仕上げの \ui{\rotl{} 左へ90°}\ui{右へ90° \rotr{}} で直してください）。
\end{caution}

% ------------------------------------------------------------------------------
\section{ダーク・フラットで補正する}\label{sec:calibration}

\begin{enumerate}[itemsep=2pt]
  \item 本番の動画を追加して選んでおきます。
  \item 品質評価タブの「入力の前処理」で \ui{ダークを選ぶ…}\ui{フラットを選ぶ…} を押し、ダーク・フラットの動画（または静止画・フォルダ）を選びます。
  \item 選んだものの名前が表示されます。あとはふつうに \ui{品質評価} から進めます。
\end{enumerate}

ダーク・フラットは全フレームを平均した「マスター」にして、色をそろえる前の画素に掛けます。
やめるときは \ui{補正をやめる} を押します。撮り方の目安は第\ref{chap:quality}章を見てください。

% ------------------------------------------------------------------------------
\section{プリセットで設定を保存する}\label{sec:preset}

\LS は起動するたびに設定が初期値に戻ります。よく使う設定の組み合わせは、\textbf{プリセット}として保存しておきましょう。

\begin{itemize}[itemsep=2pt]
  \item \textbf{保存}：右上の \ui{プリセットを保存…} を押し、名前（例「木星・シーイング普通」）を付けます。
        品質評価からスタック・仕上げ・書き出しまで、\textbf{すべてのタブの設定}がまとめて保存されます。
  \item \textbf{読み込み}：\ui{プリセット…} のメニューから名前を選ぶと、その設定に切り替わります。
  \item 処理範囲・切り抜く枠・手で置いた位置合わせ領域は、画像ごとのものなので保存されません。
  \item 保存場所は \texttt{\textasciitilde/Library/Application Support/LunaStack/Presets/} です（ほかの Mac へコピーして使えます）。
\end{itemize}

% ------------------------------------------------------------------------------
\section{途中から再開する（サイドカー）}\label{sec:sidecar}

品質評価とアライメントの結果は、元のファイルの隣に自動で保存されます（\texttt{jupiter.ser.lstkq}・\texttt{jupiter.ser.lstk} のような名前）。
次に同じファイルを開くと、この結果を読み込むので、\textbf{すぐに \ui{スタック} から始められます}。
採用率を変えた再スタックが速いのもこの仕組みのおかげです。

\begin{itemize}[itemsep=2pt]
  \item アライメントの条件を変えると、自動でやり直しが必要になります（古い結果を黙って使うことはありません）。
  \item サイドカーを消しても元の動画には影響しません。消したときは、品質評価からやり直すだけです。
  \item 入力の元のファイルを書き換えることはありません。
\end{itemize}

% ------------------------------------------------------------------------------
\section{処理を中断する・終了する}

処理中は、画面右下の \ui{中断}（\key{⌘}\key{.}）で止められます。処理中でもプレビューの拡大や表示の切り替えはできます。
処理中にウインドウを閉じたり \LS を終了したりすると、確認してから中断します。
長い処理が終わると、macOS の通知でお知らせします。
